% ===========================================================================
% Poste 17 — Sécurité de la vieillesse (PSV + SRG + Allocation)
% ===========================================================================
\poste{17}{Sécurité de la vieillesse (PSV, SRG et Allocation)}{CA\_psv}

\pdesc Le pilier public fédéral du revenu de retraite : la \emph{pension de
la Sécurité de la vieillesse} (PSV) pour les 65~ans et plus, le
\emph{Supplément de revenu garanti} (SRG) pour les aînés à faible revenu,
un \emph{supplément complémentaire} du SRG (\emph{top-up}) pour les très
faibles revenus, et l'\emph{Allocation} pour le conjoint de 60 à 64~ans
d'un prestataire du SRG.

\pobj Garantir un revenu de base à la retraite (PSV, quasi universelle) et
protéger les aînés à faible revenu (SRG, Allocation), avec récupération
auprès des plus fortunés.

\pregle Les montants officiels sont révisés \emph{chaque trimestre} ; les
paramètres annuels du modèle sont la moyenne des quatre barèmes
trimestriels --- d'où leurs décimales. Quatre composantes :

\begin{enumerate}
  \item \emph{PSV} (par adulte de 65~ans et plus, sur son revenu
        individuel $r$) : pension de base $O$, majorée de \pc{10} à
        75~ans, moins l'impôt de récupération de \pc{15} du revenu
        (pension comprise) excédant le seuil $S_r$ :
        \[
        PSV = \pos{O' - 0{,}15 \times \pos{r + O' - S_r}},
        \qquad O' = O + \mathrm{(75~ans~et~plus)} \times O_{75}
        \]
  \item \emph{SRG} (65~ans et plus) : montant maximal réduit selon le
        revenu de retraite \emph{combiné} du ménage (excluant la PSV), par
        tranches entières de 24~\$ (personne seule, taux \pc{50}) ou de
        48~\$ (couple, taux \pc{25} par conjoint) :
        \[
        SRG_{seul} = \pos{G_{seul} - 0{,}50 \times
        \Bigl\lfloor \tfrac{R}{24} \Bigr\rfloor \times 24}
        \qquad
        SRG_{couple} = \pos{G_{couple} - 0{,}25 \times
        \Bigl\lfloor \tfrac{R}{48} \Bigr\rfloor \times 48}
        \;\;\text{(par conjoint)}
        \]
  \item \emph{Supplément complémentaire} (top-up) : montant additionnel
        $T$, réduit de \pc{25} du revenu excédant une exemption $E_t$,
        par tranches arrondies \emph{au plafond} (48~\$ seul, 96~\$
        couple) :
        \[
        topup = \pos{T - 0{,}25 \times
        \Bigl\lceil \tfrac{\pos{R - E_t}}{t} \Bigr\rceil \times t}
        \]
  \item \emph{Allocation} (conjoint de 60--64~ans dont l'autre a 65~ans et
        plus) : montant maximal $A$ réduit à deux taux --- \pc{75} par
        tranches de 48~\$ jusqu'au seuil $S_a$ (la portion qui représente
        la PSV du conjoint), puis \pc{25} par tranches de 24~\$ au-delà.
\end{enumerate}

\emph{Couple mixte (un conjoint de 65~ans et plus, l'autre de 60 à
64~ans).} Le conjoint aîné reçoit la PSV et un SRG calculé sur le revenu
au-delà du seuil de l'Allocation (arrondi à la tranche supérieure de
48~\$) ; le conjoint cadet reçoit l'Allocation. Si le cadet a moins de
60~ans, il ne reçoit rien (la PSV de l'aîné demeure).

\emph{Traitement fiscal.} Seule la \emph{pension} (PSV) est imposable. Le
SRG, le supplément complémentaire et l'Allocation ne le sont pas --- mais
ils doivent être \emph{inclus dans le revenu net} qui module les crédits
d'impôt et la prime RAMQ (encadré de la section~\ref{sec:bases}). Le moteur
expose pour cela la part non imposable de chaque adulte.

\pparams

\begin{center}
\small
\begin{tabular}{l r r}
\toprule
Paramètre (moyennes annuelles des barèmes trimestriels) & 2025 & 2026 \\
\midrule
Pension de base, 65--74 ans ($O$)        & \D{8791.14}  & \D{8966.96} \\
Supplément de pension, 75 ans et + ($O_{75}$) & \D{879.15} & \D{896.70} \\
Seuil de récupération de la PSV ($S_r$)  & \D{93454}    & \D{95323} \\
Taux de récupération                     & \pc{15}      & \pc{15} \\
\midrule
SRG maximal --- personne seule ($G_{seul}$)  & \D{11096.85} & \D{11318.79} \\
SRG maximal --- par conjoint ($G_{couple}$)  & \D{7327.65}  & \D{7474.20} \\
Tranche de réduction --- seul / couple       & 24 / 48 \$   & 24 / 48 \$ \\
Taux de réduction --- seul / couple          & 50 / 25 \%   & 50 / 25 \% \\
\midrule
Top-up maximal --- seul ($T$)            & \D{2033.97}  & \D{2074.65} \\
Top-up maximal --- couple                & \D{1152.60}  & \D{1175.65} \\
Exemption du top-up --- seul ($E_t$)     & \D{2047.99}  & \D{2047.99} \\
Exemption du top-up --- couple           & \D{4095.99}  & \D{4095.99} \\
Tranche du top-up --- seul / couple ($t$)& 48 / 96 \$   & 48 / 96 \$ \\
\midrule
Allocation maximale ($A$)                & \D{16695.09} & \D{17028.99} \\
Seuil entre les deux taux ($S_a$)        & \D{11760}    & \D{12000} \\
Taux de réduction 1 / 2                  & 75 / 25 \%   & 75 / 25 \% \\
\bottomrule
\end{tabular}
\end{center}

\pnotes Le calculateur officiel comporte, pour le supplément complémentaire
des couples mixtes, une \emph{asymétrie sans justification économique} : le
montant versé dépend de la position du conjoint aîné dans le formulaire
(adulte~1 ou adulte~2) et, dans un cas, de l'âge exact de l'allocataire
(facteur $\frac12\times$, $1\times$ ou $2\times$). Le moteur reproduit ce
comportement à l'identique --- documenté comme défaut de la référence ---
afin de préserver la parité au cent. Le SRG d'un couple mixte est par
ailleurs calculé sur le revenu diminué du seuil de l'Allocation \emph{moins
un cent}, autre idiosyncrasie reproduite.

% ============================================================================
% GÉNÉRÉ par scripts/exemples-guide.ts — NE PAS ÉDITER À LA MAIN.
% Régénérer : npm run exemples (chaque chiffre est vérifié contre le moteur).
% ============================================================================
\pexemples Les paramètres sont des moyennes annuelles des quatre barèmes
trimestriels (d'où les décimales). Le SRG et son supplément se calculent par
\emph{tranches} de revenu (\num{24}, \num{48} ou
\num{96}~\$) — le revenu de pension est arrondi à la tranche
avant d'appliquer le taux de réduction.

\medskip\noindent\textbf{M10 --- retraité vivant seul, 70~ans, revenu de pension privée de \D{20000}.}
Pension de base : \num{8791.14} (le revenu est loin du seuil de
récupération de \num{93454}). SRG (taux de \pc{50},
tranches de \num{24}~\$) :
\begin{align*}
\text{revenu arrondi} &= \Bigl\lfloor \tfrac{\num{20000}}{\num{24}} \Bigr\rfloor \times \num{24} = \num{19992} \\
SRG &= \pos{\num{11096.85} - \num{19992} \times \pc{50}} = \num{1100.85}
\end{align*}
Le supplément complémentaire est éteint à ce revenu (\num{0}).
Total : \num{8791.14} $+$ \num{1100.85} $=$ \D{9891.99}.

\medskip\noindent\textbf{M13 --- couple de retraités, 68 et 66~ans, pensions privées de \D{12000} et \D{6000}.}
Couple dont les deux conjoints ont 65~ans et plus, revenu combiné de
\num{18000}. Chacun reçoit la pension de base (\num{8791.14}) et le SRG au
barème « couple » (maximum \num{7327.65}, taux \pc{25},
tranches de \num{48}~\$) :
\begin{align*}
SRG &= \pos{\num{7327.65} - \num{18000} \times \pc{25}} = \num{2827.65} \quad (\text{par conjoint}) \\
\text{supplément du couple} &= \pos{\num{1152.60} - \num{13920} \times \pc{25}} = \num{0} \\
\text{total} &= 2 \times \num{8791.14} + 2 \times \num{2827.65} + \num{0} = \num{23237.58}
\end{align*}
Sécurité de la vieillesse du ménage : \D{23237.58}.

\medskip\noindent\textbf{M11 --- couple de retraités, 72 et 70~ans, pensions privées de \D{30000} et \D{10000}.}
Même structure, revenu combiné de \num{40000} : la réduction du SRG
(\num{39984} $\times$ \pc{25} $=$
\num{9996}) dépasse le maximum de
\num{7327.65} : SRG \emph{éteint}. Reste la pension de base de chacun :
$2 \times \num{8791.14} = \num{17582.28}$, soit \D{17582.28}.

\medskip\noindent\textbf{M12 --- retraité vivant seul, 76~ans, revenu de pension privée de \D{110000}.}
75~ans et plus : pension majorée (\num{8791.14} $+$ \num{879.15}
$=$ \num{9670.29}). Le revenu déclenche l'\emph{impôt de récupération} :
\begin{align*}
\text{récupération} &= \pos{\num{110000} + \num{9670.29} - \num{93454}} \times \pc{15} = \num{3932.44} \\
PSV &= \num{9670.29} - \num{3932.44} = \num{5737.85}
\end{align*}
Pension nette : \D{5737.85} (la pension s'éteint complètement vers
\D{148252.31} de revenu).


\prefs Loi sur la sécurité de la vieillesse (L.R.C. 1985, ch.~O-9) ; impôt
de récupération : Loi de l'impôt sur le revenu, art.~180.2 ; Emploi et
Développement social Canada (montants trimestriels) ; ministère des
Finances du Québec, \emph{Dépenses fiscales}, fiche 110113 (traitement du
SRG dans le revenu net).
